
\chapter{Agente de Inteligência Artificial}
\label{chap:agente-ia}

% Apresentar a finalidade do agente de IA no DoseCerta
% e como ele contribui para a interação com o usuário.

\section{Objetivo e Requisitos}
\label{sec:objetivo-agente-ia}

% Descrever o problema atendido pelo agente e suas funções previstas.
% Relacionar sua atuação aos requisitos funcionais e não funcionais.

\section{Tecnologias e Modelo Utilizado}
\label{sec:tecnologias-agente-ia}

% Informar modelo, serviço, biblioteca ou framework utilizado,
% quando aplicável. Descrever critérios de seleção.
% Não afirmar que houve treinamento próprio se não ocorreu.

\section{Arquitetura do Agente}
\label{sec:arquitetura-agente-ia}

% Descrever os componentes do agente e suas responsabilidades:
% entrada do usuário, processamento, interpretação e resposta.
% Inserir diagrama da arquitetura, se disponível.

\section{Interpretação de Linguagem Natural}
\label{sec:interpretacao-linguagem-natural}

% Explicar como as mensagens do usuário são interpretadas
% para identificar intenções e informações relevantes.
% Apresentar exemplos reais de entradas e saídas, se disponíveis.

\section{Fluxo de Interação}
\label{sec:fluxo-interacao-agente}

% Descrever o fluxo entre usuário, agente e demais módulos.
% Explicar como o agente solicita informações adicionais
% ou confirma dados, se essas funções existirem.

\section{Integração com o Backend e o Aplicativo}
\label{sec:integracao-agente-ia}

% Descrever como o agente recebe solicitações e devolve resultados.
% Informar interfaces, endpoints ou formatos de dados utilizados,
% conforme a implementação. Evitar repetir os capítulos 10 e 11.

\section{Validação das Informações e Restrições}
\label{sec:validacao-agente-ia}

% Descrever como as informações interpretadas são verificadas
% antes de serem utilizadas pelo sistema.
% Explicar limites de atuação e tratamento de respostas inválidas.
% Destacar que informações relacionadas a medicamentos exigem
% confirmação e não devem ser presumidas pelo agente.

\section{Privacidade e Segurança}
\label{sec:privacidade-agente-ia}

% Descrever o tratamento de dados enviados ao agente,
% incluindo medidas efetivamente adotadas para proteger
% informações pessoais e sensíveis.

\section{Limitações da Solução}
\label{sec:limitacoes-agente-ia}

% Registrar limitações observadas ou conhecidas:
% ambiguidades, erros de interpretação, dependência de conexão
% ou do serviço utilizado, conforme aplicável.

\section{Síntese do Agente de IA}
\label{sec:sintese-agente-ia}

% Resumir a contribuição do agente ao DoseCerta,
% as funções implementadas e o estágio atual de desenvolvimento.
